\chapter{Reportes de ventas}

\pantalla{Reportes} muestra las ventas \textbf{cobradas}, según la hora en que se cobraron. Al abrirla trae el resumen de hoy.

\section{Elegir qué ver}

\captura{reportes-resumen}{Resumen de ventas}

\begin{enumerate}
  \item Elija el tipo: \boton{Resumen de ventas} o \boton{Ventas detalladas}.
  \item Elija las fechas \textbf{Desde} y \textbf{Hasta} con el calendario, o use los atajos \boton{Hoy}, \boton{Ayer}, \boton{Esta semana} y \boton{Este mes}.
  \item Si quiere, elija un \textbf{Empleado}: trae las ventas que esa persona registró, a las que agregó un envío o que cobró. También se puede elegir el DIRECTORIO. Los empleados dados de baja aparecen marcados como \textit{de baja}.
  \item Presione \boton{Generar}.
\end{enumerate}

Arriba de la tabla están los totales del período: cantidad de ventas, total, efectivo, QR y ticket promedio.

\begin{description}
  \item[Resumen de ventas] tiene una fila por día con sus ventas, efectivo, QR y total.
  \item[Ventas detalladas] tiene una fila por venta: número, hora de cobro, mesa, quién la abrió y quién la cobró, y los montos.
\end{description}

\captura{reportes-detalladas}{Ventas detalladas}

\section{Detalle de una venta}

\captura{reportes-venta}{Detalle de una venta}

El ojo de cada fila abre la venta completa:

\begin{itemize}
  \item Cada \textbf{envío} con su hora y quién lo registró, en una tabla con \textbf{Cant.}, \textbf{Producto}, \textbf{Consumo} (Local o Para llevar) y \textbf{Subtotal}. Lo quitado se lee en rojo debajo del producto.
  \item El \textbf{cobro}: efectivo, QR, total, recibido, cambio, quién cobró y cuándo.
  \item Cuántas veces se imprimió el ticket.
\end{itemize}

\boton{Reimprimir ticket} vuelve a imprimirlo, con el mismo límite de impresiones de Caja (solo quien tiene acceso a Caja).

\section{Descargar o enviar por correo}

\begin{itemize}
  \item \boton{Descargar PDF} guarda el reporte que se está viendo. El PDF de ventas detalladas trae cada venta con su tabla de productos.
  \item Para mandarlo, escriba un correo en \textbf{Enviar a este correo} y presione \boton{Enviar por correo}: llega con el PDF adjunto.
\end{itemize}

\begin{importante}
Para que los correos salgan de verdad, la computadora debe tener configurada la clave del servicio de correo (Brevo), la misma que usa la recuperación de contraseñas. Pídasela al encargado del sistema.
\end{importante}

\section{Envíos programados (solo el DIRECTORIO)}

El DIRECTORIO puede hacer que el reporte llegue solo por correo ciertos días a cierta hora.

\captura[0.8]{reportes-nuevo-envio}{Nuevo envío programado}

\begin{enumerate}
  \item Presione \boton{Envíos programados} y luego \boton{Nuevo envío}.
  \item Escriba el \textbf{correo de destino}, elija el \textbf{reporte}, marque los \textbf{días de la semana} y la \textbf{hora}.
  \item Presione \boton{Guardar}.
\end{enumerate}

\captura[0.8]{reportes-programados}{Lista de envíos programados}

En la lista, cada envío muestra cuándo salió el último y cuándo sale el próximo. Con sus botones se puede:

\begin{itemize}
  \item \textbf{Pausar} o reanudar con el interruptor.
  \item \textbf{Enviar ahora}, sin esperar el día y la hora.
  \item \textbf{Modificar} con el lápiz o \textbf{borrar} con el tacho.
\end{itemize}

Cómo funciona:

\begin{itemize}
  \item Cada envío trae las ventas \textbf{desde el envío anterior} hasta ese momento.
  \item Si la computadora estaba apagada a esa hora, el reporte sale apenas se enciende el sistema y en la lista figura como \textit{Enviado al encender}.
  \item Si el correo falla, el sistema lo vuelve a intentar solo a los 15 minutos.
\end{itemize}
